We reframed trip-scale route generation around a distinction between \emph{structural realism} (geometric validity) and \emph{behavioral realism} (contextual plausibility). The diagnostic finding---that map-free continuous generation fails to converge to road corridors even with soft priors---motivates treating corridor choice as a discrete decision on the road graph rather than a continuous optimization.

CascadeTraj operationalizes this insight: corridor selection on a road graph provides structural validity, while temporal--spatial semantic conditioning provides behavioral grounding. The key contribution is not the architecture per se, but the framing: behavioral realism means that the same OD under different contexts should induce different corridor distributions. Temporal--spatial semantics are not auxiliary features---they are the mechanism that turns diversity into meaningful behavioral variation.

\paragraph{Limitations and future directions.}
The approach depends on road-graph quality and requires robust corridor labeling from trajectories. Extending semantic features to time-dynamic POI activity (e.g., open hours) and evaluating cross-city transfer under consistent corridor definitions are natural next steps.
